\section{Case Studies}
{\color{sub}\footnotesize 2 hrs \gap$\cdot$\gap asked in \mk{14 of 29 papers}
\gap$\cdot$\gap 5 syllabus topics \gap$\cdot$\gap \mk{92 marks, 3.9\% of the archive}
$\rightarrow$ small, but a near-certain 4-mark short note since 2079.}

\vspace{3pt}\noindent
\colorbox{gdL}{\makebox[\dimexpr\textwidth-2\fboxsep][l]{%
  \textcolor{gd}{\bfseries\footnotesize READ THIS FIRST}}}
\par\nobreak\vspace{3pt}

Only two things are ever asked, in rotation: \mk{objectives of case study} (8 papers) and
\mk{steps / phases of case study} (5 papers, all since 2079). Learn the \textbf{6 steps}
and the \textbf{5 phases} as two separate numbered lists; examiners use both words. Twice a
paper has set an applied 8-marker (prepare a case study of NEA or Nepal Telecom; recommend
improvements for your college): answer it by filling in the \textbf{report structure} of
$\S$4.3.

Tier chips count \textbf{papers out of 29}: \tS{n} 8+ \gap \tF{n} 4--7 \gap \tP{n} 1--3
inside an 8-mark question. Bold year = Regular exam.

\hr

%=======================================================================
\T{4.1 Introduction and Objectives}

\Q{What is a case study? Why are case studies conducted? Objectives and value of case study}{\m{2}\m{3}\m{4}}
{\color{sub}\footnotesize \tS{8} \yr{81 Ba, 80 Ba, \textbf{73 Ch}, \textbf{71 Ch}, 71 Shr, 70 Asa, \textbf{69 Ch}, \textbf{66 Bh}} \gap$\cdot$\gap short note in 71 Shr, 69 Ch, 66 Bh (``value of case study'')}

\sbs{%
\lead{Meaning}
\begin{itemize}
  \item A \textbf{case} = written description of a \textbf{real management situation}:
    facts, reports, even opinions, but \textbf{no evaluation}.
  \item \textbf{Case study method}: learners analyse the case, identify problems, weigh
    alternatives and \textbf{decide}, as a manager would.
  \item Belief behind it: managerial competence comes from studying and discussing
    \textbf{actual situations}, not only theory.
  \item Also a \textbf{research method}: in-depth study of one unit (firm, project, event) to
    understand causes or test a hypothesis.
  \item Harvard Business School made it the standard method of management teaching
    (1920s) \added.
\end{itemize}
\lead{Why conducted (AJ Sir)}
\begin{itemize}
  \item Teach managerial skills: \textbf{analyse} a situation, \textbf{frame} the right
    questions, \textbf{find answers}, \textbf{make decisions}.
  \item Link the three foundations of management practice: \textbf{research, theory,
    practice}.
\end{itemize}}{%
\lead{Objectives}
\begin{enumerate}
  \item \textbf{Impart knowledge and facts} about real organizations.
  \item Develop \textbf{problem analysis} skill: separate symptoms from causes.
  \item Develop \textbf{decision-making} skill under incomplete information.
  \item Improve \textbf{communication}: discuss, defend, write reports.
  \item Shape \textbf{attitude}: nothing in human behaviour is absolutely right or wrong;
    respect other views.
  \item Apply \textbf{theory to practice}; test principles against reality.
  \item Build \textbf{teamwork} and judgement through group discussion.
\end{enumerate}
\lead{Value (short note, 66 Bh)}
\begin{itemize}
  \item $+$ realistic, active learning; develops analytical thinking; learner participates;
    many solutions welcome; bridges classroom and industry.
  \item $-$ time consuming (gathering, arranging, writing facts is long and tedious); needs
    a skilled discussion leader; one case may not generalize.
\end{itemize}}

\creamq{Types of case studies}{\m{2}}
{\color{sub}\footnotesize \tP{1} \yr{80 Ba}}
\sbs{%
\begin{itemize}
  \item \textbf{Illustrative}: descriptive; one or two instances show what a situation is
    like; makes the unfamiliar familiar.
  \item \textbf{Exploratory (pilot)}: short study before a large investigation, to find the
    right questions and measures. Risk: early findings released as conclusions.
\end{itemize}}{%
\begin{itemize}
  \item \textbf{Cumulative}: combines past studies from several sites and times $\rightarrow$
    wider generalization at no new cost.
  \item \textbf{Critical instance}: one or a few sites of unique interest, or to challenge a
    universal claim; answers cause-effect questions.
\end{itemize}}
{\small Teaching cases may also be classed as \textbf{incident} (short event), \textbf{problem /
decision} (choose an action), \textbf{appraisal} (evaluate a situation) and \textbf{descriptive}
cases.}

\penalty0\vspace{2pt}

%=======================================================================
\T{4.2 Phases and Steps of Case Study}

\Q{Explain the phases and steps of case study, with an example from your field}{\m{4}\m{5}\m{8}}
{\color{sub}\footnotesize \tF{5} \yr{\textbf{82 Bh}, \textbf{81 Bh}, 81 Ba, \textbf{80 Bh}, \textbf{79 Bh}} \gap$\cdot$\gap short note in 82 Bh, 81 Bh \gap$\cdot$\gap \mk{every paper since 2079 except 80 Ba}}

\sbs{%
\lead{Five phases (group discussion)}
\begin{enumerate}
  \item \textbf{Study the details}: each member reads the case. What is happening? What
    more data is needed?
  \item \textbf{Collect additional facts}: ask the discussion leader; organize facts by
    relevance to the decision.
  \item \textbf{Define the problem}: the group agrees on the real problem and phrases the
    critical questions.
  \item \textbf{Individual decisions}: each writes ``How would I handle it? How would I
    support my decision?''; members with the same view form sub-groups and argue their case.
  \item \textbf{Learn from the whole case}: ``How could more have been achieved?''
\end{enumerate}}{%
\lead{Six steps (individual analysis)}
\begin{enumerate}
  \item \textbf{State the problem}: what is the case about, what must be done or answered.
  \item \textbf{Collect and analyse data}: arrange facts systematically; separate facts from
    opinions.
  \item \textbf{Formulate tentative solutions}: list all alternatives, keep an open mind.
  \item \textbf{Select the recommended solution}: best answer to the stated problem on the
    given data.
  \item \textbf{Prepare the written report}.
  \item \textbf{Checklist}: what does the solution achieve? what difficulties must be
    overcome? what limits (time, money, equipment, manpower)? is it the best alternative?
\end{enumerate}}

\lead{Example from engineering: why did a hydropower project finish two years late? \added}
{\small
\begin{tabularx}{\textwidth}{@{}L{3.1cm}X@{}}
\toprule
\hrow \thead{Step} & \thead{Applied} \\
\midrule
1 State the problem & 60 MW project, planned 4 years, took 6; cost overrun 40\%. Why, and how to prevent it next time? \\
2 Collect and analyse data & Schedule vs actual (Gantt), change orders, land-acquisition dates, contractor payments, site logs, interviews with engineers and locals. Finding: land dispute (8 months), tunnel geology surprise (10 months), late payments $\rightarrow$ contractor slowdown. \\
3 Tentative solutions & (a) better geological survey; (b) settle land before tender; (c) escrow payment to contractor; (d) matrix project organization with one accountable manager. \\
4 Recommended solution & Combine (a)+(b) as pre-construction gates, (c) for cash flow; (d) as the structure. \\
5 Written report & Background, problem, data, analysis, recommendations, implementation plan. \\
6 Checklist & Achieves on-time completion; needs upfront survey budget; limit: political land issues; still the best alternative. \\
\bottomrule
\end{tabularx}}

\penalty0\vspace{2pt}

%=======================================================================
\T{4.3 Case Report Writing, and Preparing a Case Study}

\Q{Prepare a case study of an organization (NEA / Nepal Telecom; your college) following the case structure}{\m{8}}
{\color{sub}\footnotesize \tP{2} \yr{\textbf{74 Ch}, \textbf{72 Ch}} \gap$\cdot$\gap applied 8-markers: fill in the structure}

\sbs{%
\lead{Case report writing methodology}
\begin{enumerate}
  \item \textbf{Organizational background}: goals, formation, structure, social value,
    history, financial position.
  \item \textbf{Identification of the problem}.
  \item \textbf{Hypothesis}: the suspected cause.
  \item \textbf{Case research design}: what data, from whom, how.
  \item \textbf{Collection and analysis of data}.
  \item \textbf{Generalization and interpretation}.
  \item \textbf{Conclusions and recommendations}.
\end{enumerate}
}{%
\lead{Typical report titles (a full O\&M case study)}
{\small Acknowledgement; organization (history, objectives, input-process-output); structure
(chart); form of ownership; personnel policies; manpower planning; recruitment and selection;
training; job evaluation; merit rating; wages and incentives; motivation; marketing;
leadership; MIS; drawbacks; suggestions; conclusion; references.}}

\lead{NEA or Nepal Telecom (72 Ch): planning horizon, leadership, motivation, HRD}
\begin{itemize}
  \item \textbf{Background}: NEA, government-owned (NEA Act 2041), generation, transmission,
    distribution; long load-shedding (up to 18 h/day in 2015--16).
  \item \textbf{Problem}: poor performance: losses, leakage (about 25\% system loss),
    load-shedding, low morale.
  \item \textbf{Hypothesis}: weak leadership and short planning horizon, not only lack of
    power.
  \item \textbf{Data}: supply-demand, loss figures, staff surveys, interviews.
  \item \textbf{Planning horizon}: short term (cut leakage, import power, manage
    dispatch); medium (transmission lines); long (new generation, export).
  \item \textbf{Leadership}: decisive, visible, results-driven MD ended load-shedding in
    2016--18 through leakage control and better dispatch: a lesson in leadership.
  \item \textbf{Motivation}: performance targets per distribution centre, recognition,
    incentives, protection from political interference.
  \item \textbf{HRD}: technical training, merit-based promotion, succession planning.
  \item \textbf{Recommendations} + implementation plan and monitoring.
\end{itemize}

\creamq{Recommendations to your college management (74 Ch)}{\m{8}}
\sbs{%
\begin{itemize}
  \item Frame it as a case: \textbf{background} (programmes, students, faculty),
    \textbf{problems} observed over 3+ years, \textbf{analysis}, \textbf{recommendations}.
  \item \textbf{Academic}: regular classes and full-time faculty; faculty training;
    updated labs; industry projects and internships; continuous assessment.
  \item \textbf{Infrastructure}: library e-resources, internet, classrooms, hostels.
\end{itemize}}{%
\begin{itemize}
  \item \textbf{Student support}: counselling, career cell, alumni network, clubs.
  \item \textbf{Research}: research grants, publication support, MoUs with industry.
  \item \textbf{Management}: MIS for attendance and results, student feedback, transparent
    fees, accreditation (UGC QAA).
  \item \textbf{Image}: website, social media, results publicity, community programmes.
  \item Close with priorities and a time-bound plan.
\end{itemize}}

\penalty0\vspace{2pt}

%=======================================================================
\T{4.4 PYQ Mapping}

\lead{Theory questions, covered in this document}
\begin{itemize}
  \item \textbf{What is a case study, why conducted, objectives, value} $\rightarrow$
    $\S$4.1. 8 papers; three of them as a 4-mark short note.
  \item \textbf{Types of case studies} \yr{80 Ba \m{2}} $\rightarrow$ $\S$4.1.
  \item \textbf{Phases and steps, with an example from your field} $\rightarrow$ $\S$4.2.
    \yr{\textbf{82 Bh} \m{4}, \textbf{81 Bh} \m{4}, 81 Ba \m{5}, \textbf{80 Bh} \m{8}, \textbf{79 Bh} \m{8}}
  \item \textbf{Prepare a case study of NEA / Nepal Telecom; recommendations for your
    college} \yr{\textbf{74 Ch} \m{8}, \textbf{72 Ch} \m{8}} $\rightarrow$ $\S$4.3.
  \item 73 Ch and 71 Ch pair case-study objectives with MIS needs; the MIS half is in Ch.\,5.
\end{itemize}

\lead{Numerical content}
\begin{itemize}
  \item \textbf{This chapter has no numerical component.}
\end{itemize}

\lead{Sources}
\begin{itemize}
  \item Adhikari pp259--272 and AJ Sir's 7-slide deck cover all five syllabus topics
    between them. S.K.\ Joshi's textbook follows an older syllabus and has no case-study
    chapter, so nothing was taken from it here.
\end{itemize}
